\chapter{Rilevazione e correzione degli errori}

Il ricevente sa ormai trovare dove un frame inizia e dove finisce, ed è tutto ciò che sa fare. Con il solo payload non c'è nulla con
cui confrontarlo: ogni sequenza di bit è una sequenza di bit legittima. Il mittente calcola quindi qualcosa di \textbf{ridondante}
sui dati e lo invia insieme ad essi, nel trailer in fondo al frame; il ricevente ripete il calcolo su ciò che è arrivato e confronta.
In questo capitolo vediamo che cosa sia questo qualcosa, quanto sia forte, e che cosa succede quando non lo è abbastanza.

% ══════════════════════════════════════════════════════════════
\section{La rilevazione degli errori}
% ══════════════════════════════════════════════════════════════

\dfn{Il problema}{
    Siano $D$ i nostri dati, di $d$ bit: non solo il datagramma, ma anche l'\textbf{header di collegamento} (indirizzi, numeri di
    sequenza, ogni campo). Per proteggerli dagli errori sui bit aggiungiamo al messaggio dei \textbf{bit di rilevazione e correzione
    degli errori}, \textbf{EDC} (\emph{Error Detection and Correction}), calcolati su $D$ prima che parta. Arrivano $D'$ ed $EDC'$,
    che possono differire da entrambi: il ricevente ricalcola l'EDC su $D'$ e lo confronta con $EDC'$. L'obiettivo è capire se
    qualcosa è cambiato e, se possibile, ricostruire gli originali.
}

\begin{figure}[H]
\centering
\begin{tikzpicture}[every node/.style={font=\scriptsize\sffamily}]
  \node[draw, fill=lnet!50, minimum width=1.4cm, minimum height=0.6cm] (dg) at (0,1) {datagramma};
  \node[draw, fill=llink!60, minimum width=1.2cm, minimum height=0.6cm, right=0pt of dg] {header link};
  \node[draw, fill=cloudbg, minimum width=2.6cm, minimum height=0.6cm] (d) at (0.6,0) {D};
  \node[draw, fill=netred!25, minimum width=0.8cm, minimum height=0.6cm, right=0pt of d] (e) {EDC};
  \node[draw, fill=cloudbg, minimum width=2.6cm, minimum height=0.6cm] (d2) at (9.2,0) {D'};
  \node[draw, fill=netred!25, minimum width=0.8cm, minimum height=0.6cm, right=0pt of d2] (e2) {EDC'};
  \draw[very thick, decorate, decoration={snake, amplitude=1.5pt, segment length=8pt}, -{Stealth}] (e.east) -- node[above=3pt] {collegamento soggetto a errori} (d2.west);
  \node[diamond, draw, aspect=2.5, fill=white, inner sep=1pt] (q) at (9.6,1.3) {D'+EDC' = D+EDC ?};
  \draw[-{Stealth}] (d2.north) -- (q);
\end{tikzpicture}
\caption{Il ricevente usa EDC' per verificare se D' è integro.}
\end{figure}

\warning{Rilevato, non avvenuto}{
    Il controllo va a buon fine: vuol dire che nessun bit si è invertito? \textbf{No}: vuol dire che nessun errore è stato
    \textbf{rilevato}. Alcune combinazioni di inversioni lasciano i bit ridondanti perfettamente coerenti con i dati danneggiati. La
    domanda a cui risponde un ricevente è sempre la prima, e le parole contano: gli errori si \textbf{rilevano}, non si osservano.
    Uno schema più forte rende la probabilità più piccola, nessuno la rende zero; e la forza si paga: più bit da calcolare, più bit
    da trasmettere, più circuiti per calcolarli.
}

\subsection{Due strategie per lo stesso problema}

Entrambe aggiungono ridondanza; differiscono in \textbf{quanta}.

\begin{itemize}
    \item I \textbf{codici a rilevazione d'errore} ne aggiungono abbastanza per sapere che qualcosa è sbagliato, e chiedere di nuovo.
    \item I \textbf{codici a correzione d'errore} ne aggiungono abbastanza per ricostruire che cosa doveva essere stato inviato, e
          ripararlo sul posto. Si parla di \textbf{FEC} (\emph{Forward Error Correction}): \emph{forward} perché nulla torna
          indietro.
\end{itemize}

La scelta non riguarda quale sia più forte: riguarda \textbf{quanto costa una ritrasmissione} su quel particolare collegamento.

\begin{itemize}
    \item \textbf{Fibra, coassiale, doppino}: gli errori sono rari. Rilevarli e rimandare l'occasionale frame danneggiato costa meno
          che portare bit di riparazione su ogni frame.
    \item \textbf{Senza fili, vecchi doppini dell'ultimo miglio}: gli errori sono la norma, e una ritrasmissione ha la stessa
          probabilità di essere danneggiata del primo tentativo.
    \item Quindi su un canale rumoroso si ripara localmente, su uno pulito si chiede di nuovo: entrambe le strategie sono in uso, su
          collegamenti a un metro di distanza.
    \item Nessuna delle due gestisce tutto: il canale non tratta i bit di controllo diversamente dai bit di dati, per lui sono solo
          bit. Il codice deve essere abbastanza forte per gli errori \textbf{attesi su quel collegamento}.
\end{itemize}

% ══════════════════════════════════════════════════════════════
\section{Il controllo di parità}
% ══════════════════════════════════════════════════════════════

\dfn{Controllo di parità}{
    Il \textbf{controllo di parità} è la forma più semplice di rilevazione: si aggiunge al messaggio \textbf{un solo bit di parità},
    scelto in modo che il numero totale di 1 nei $d+1$ bit sia \textbf{pari} (schema a parità pari) o \textbf{dispari} (schema a parità
    dispari).
}

\begin{table}[H]
\centering
\begin{tabular}{lcc}
\toprule
\textbf{Messaggio} & \textbf{Bit di parità (pari)} & \textbf{Bit di parità (dispari)} \\
\midrule
\texttt{10101100} (quattro 1) & 0 & 1 \\
\texttt{11010000} (tre 1) & 1 & 0 \\
\bottomrule
\end{tabular}
\caption{Bit di parità pari e dispari.}
\end{table}

Equivalentemente, con la parità pari il bit di parità è lo \textbf{XOR di tutti i bit di dati}: un solo bit di overhead, per
quanto lungo sia il messaggio. Il ricevente deve solo contare gli 1 nei $d+1$ bit ricevuti: se con uno schema pari trova un numero
dispari di 1 (o viceversa), sa che si è verificato \textbf{almeno un errore}. Non gli serve memoria di ciò che è stato inviato, né una
seconda copia con cui confrontare: conta, e il conteggio o rispetta la convenzione o no. Per questo basta un bit di hardware, alla
velocità a cui arrivano i bit.

\begin{esempio}{Un errore rilevato e uno no}
Si invia \texttt{10101100} con parità pari: \texttt{10101100 0}.
\begin{itemize}
    \item Si riceve \texttt{1010\textcolor{netred}{0}100 0}: tre 1, numero dispari $\Rightarrow$ \textbf{errore rilevato}.
    \item Si riceve \texttt{1010\textcolor{netred}{0}10\textcolor{netred}{1} 0}: quattro 1, numero pari $\Rightarrow$ due errori, \textbf{non rilevati}.
\end{itemize}
\end{esempio}

Se si invertono due bit dello stesso frame, il numero di 1 cambia di due e la sua parità non cambia affatto: il conteggio torna,
il frame è accettato e il danno passa al livello superiore. La parità singola rileva un numero \textbf{dispari} di inversioni, e
nient'altro: due, quattro, sei sono invisibili. Nel linguaggio di fine capitolo, è un codice di \textbf{distanza due}.

\warning{Gli errori arrivano a raffica}{
    Quanto è probabile un numero pari di errori? Se le inversioni fossero rare e indipendenti, più di una nello stesso frame sarebbe
    così improbabile che un bit di parità basterebbe. Ma nelle reti non è così: gli errori sono spesso \textbf{raggruppati a raffica}
    (\emph{burst}), cioè sequenze in cui il primo e l'ultimo bit sono sbagliati e quelli in mezzo sono un misto di bit sbagliati e
    corretti. La causa fisica dura più di un tempo di bit: un affievolimento su un canale radio, un transitorio su una linea in rame,
    un graffio su un disco. In condizioni di burst la probabilità che la parità singola non veda il danno si avvicina al
    \textbf{50\%}: non è davvero utile.
}

C'è anche un rovescio della medaglia. Con blocchi di 1000 bit e un errore ogni mille bit, errori indipendenti danneggiano
\textbf{la maggior parte} dei blocchi; raffiche di cento errori ne danneggiano \textbf{uno su cento}.

\begin{esempio}{Eppure un bit di parità costa poco}
Prendiamo un cavo abbastanza buono da invertire un bit su un milione, e blocchi di mille bit.
\begin{itemize}
    \item Per \textbf{correggere} un errore singolo servono circa dieci bit di controllo per blocco: \textbf{diecimila} bit per
          megabit di dati.
    \item Per \textbf{accorgersene} basta un bit per blocco, mille per megabit; un blocco su mille risulta sbagliato e viene
          rimandato, al costo di 1001 bit. In tutto circa \textbf{duemila} bit.
\end{itemize}
Duemila contro diecimila: su un collegamento buono rilevazione e ritrasmissione vincono. Su un collegamento cattivo l'aritmetica
si rovescia, perché le ritrasmissioni si moltiplicano e ciascuna può fallire a sua volta.
\end{esempio}

\subsection{Calcolare la parità in un ordine diverso}

L'idea è calcolare i bit di controllo sui dati in un ordine \textbf{diverso} da quello in cui i bit escono sulla linea. Si
dispongono i dati in righe, si invia riga dopo riga, ma si calcola una parità \textbf{per colonna}.

\begin{figure}[H]
\centering
\begin{tikzpicture}[every node/.style={font=\small\ttfamily}, c/.style={draw, minimum size=0.5cm, inner sep=0pt}]
  \foreach \r [count=\i] in {{1,0,0,1,1,1,0},{1,1,0,0,1,0,1},{1,1,1,0,1,0,0},{1,1,1,0,1,1,1}} {
    \node[font=\scriptsize\sffamily, anchor=east] at (0.15,-\i*0.5) {riga \i};
    \foreach \v [count=\j] in \r \node[c] (c\i\j) at (\j*0.5+0.3,-\i*0.5) {\v};
  }
  \foreach \v [count=\j] in {0,1,0,1,0,0,0} \node[c, fill=llink!60] at (\j*0.5+0.3,-2.75) {\v};
  \node[font=\scriptsize\sffamily, anchor=east] at (0.15,-2.75) {parità};
  \foreach \j in {1,2,5,6,7} \node[text=netred, font=\small\bfseries\sffamily] at (\j*0.5+0.3,-3.3) {$\times$};
  \draw[netred, very thick, rounded corners=2pt] (2.55,-1.27) rectangle (4.05,-0.73);
  \draw[netred, very thick, rounded corners=2pt] (0.55,-1.77) rectangle (1.55,-1.23);
  \node[font=\scriptsize\sffamily, text width=5.2cm, align=left, anchor=west] at (4.6,-1.4) {si invia riga dopo riga: una raffica di cinque bit consecutivi (in rosso, a cavallo tra la riga 2 e la riga 3) tocca un bit in ciascuna di cinque colonne diverse, e tutte e cinque le parità di colonna falliscono};
\end{tikzpicture}
\caption{La parità per colonne cattura una raffica che attraversa le righe.}
\end{figure}

Una raffica che copre fino a $n$ bit trasmessi (con $n$ colonne) tocca \textbf{al più un bit per colonna}, e ogni colonna toccata fallisce
la propria parità. La parità non è diventata più forte: lo è diventata la \textbf{disposizione}. Un codice pensato per errori isolati
ora cattura le raffiche (la tecnica si chiama \textbf{interleaving}).

\subsection{La parità bidimensionale}

\dfn{Parità bidimensionale}{
    Nella \textbf{parità bidimensionale} i $d$ bit del messaggio vengono disposti in $i$ righe e $j$ colonne, e a \textbf{ogni riga} e
    \textbf{ogni colonna} si associa un bit di parità: $i + j + 1$ bit di controllo in tutto (l'ultimo, all'angolo, è la parità delle
    parità).
}

\begin{figure}[H]
\centering
\begin{tikzpicture}[every node/.style={font=\small\ttfamily}, c/.style={draw, minimum size=0.5cm, inner sep=0pt}]
  \def\tabella#1#2#3{
    \begin{scope}[shift={#1}]
      \foreach \r [count=\i] in {#2} {
        \foreach \v [count=\j] in \r {
          \pgfmathsetmacro{\fillc}{(\i==5 || \j==6) ? "llink!60" : "white"}
          \node[c, fill=\fillc] (c\i\j) at (\j*0.5,-\i*0.5) {\v};
        }
      }
      #3
    \end{scope}
  }
  \node[font=\footnotesize\bfseries\sffamily] at (1.75,0.2) {senza errori};
  \tabella{(0,0)}{{1,0,1,0,1,1},{1,1,1,1,0,0},{0,1,1,1,0,1},{0,0,1,0,1,0},{0,0,0,0,0,0}}{}
  \node[font=\footnotesize\bfseries\sffamily] at (5.75,0.2) {un errore};
  \tabella{(4,0)}{{1,0,1,0,1,1},{1,0,1,1,0,0},{0,1,1,1,0,1},{0,0,1,0,1,0},{0,0,0,0,0,0}}{
     \draw[netred, very thick] (0.25,-1.25) rectangle (3.25,-0.75);
     \draw[netred, very thick] (0.75,-0.25) rectangle (1.25,-2.75);
     \node[circle, draw=netred, very thick, minimum size=0.55cm] at (1,-1) {};
  }
  \node[font=\footnotesize\bfseries\sffamily] at (9.75,0.2) {due errori};
  \tabella{(8,0)}{{1,0,1,0,1,1},{1,0,1,1,1,0},{0,1,1,1,0,1},{0,0,1,0,1,0},{0,0,0,0,0,0}}{
     \draw[netred, very thick] (0.75,-0.25) rectangle (1.25,-2.75);
     \draw[netred, very thick] (2.25,-0.25) rectangle (2.75,-2.75);
  }
  \node[font=\scriptsize\sffamily, text width=3cm, align=center] at (1.75,-3.2) {in giallo i bit di parità (pari)};
  \node[font=\scriptsize\sffamily, text width=3.4cm, align=center] at (5.75,-3.2) {riga 2 e colonna 2 sbagliate: il bit all'incrocio è quello errato e si può correggere};
  \node[font=\scriptsize\sffamily, text width=3.4cm, align=center] at (9.75,-3.2) {due colonne sbagliate, righe corrette: errore rilevato ma non correggibile};
\end{tikzpicture}
\caption{Parità bidimensionale: dati $4 \times 5$ bit, con parità di riga (ultima colonna) e di colonna (ultima riga).}
\end{figure}

\begin{itemize}
    \item Un bit invertito rompe \textbf{esattamente una riga} ed \textbf{esattamente una colonna}, e l'incrocio delle due nomina il bit
          sbagliato. Nominarlo basta per rimetterlo a posto: la rilevazione è diventata \textbf{correzione}, senza rimandare indietro
          alcun messaggio.
    \item Se invece si invertono due bit della \textbf{stessa riga}, la parità di riga non cambia (le due inversioni si annullano) e
          protestano solo le due colonne: si sa che il frame è danneggiato, ma nessun bit può essere nominato.
    \item La parità bidimensionale \textbf{rileva} quindi qualsiasi combinazione di \textbf{due errori}, ma ne \textbf{corregge} uno
          solo. La distanza tra i due verbi è l'argomento della fine del capitolo.
\end{itemize}

\subsection{La correzione in avanti}

Riparare al ricevente ha un nome, \textbf{FEC}: il frame viene usato così come è arrivato, e nulla viene richiesto di nuovo. Il
guadagno non sono i bit risparmiati, ma il \textbf{round trip non atteso}, e su alcuni collegamenti quell'attesa è tutto il problema:

\begin{itemize}
    \item audio e video in tempo reale, dove un dato in ritardo è peggio di un dato sbagliato;
    \item collegamenti con lo spazio profondo, dove un round trip si misura in \textbf{minuti};
    \item dischi ottici, dove il danno è un graffio e non c'è nessuno a cui chiedere una seconda copia.
\end{itemize}

La FEC si combina anche con la ritrasmissione: si ripara ciò che si può riparare, e si chiede di nuovo il resto.

% ══════════════════════════════════════════════════════════════
\section{I checksum}
% ══════════════════════════════════════════════════════════════

La seconda famiglia di tecniche legge i dati come una sequenza di interi di $k$ bit e li \textbf{somma}. Ne abbiamo già incontrata una
al livello di trasporto: parole di 16 bit, complemento a uno della somma, nell'header del segmento.

\begin{itemize}
    \item È \textbf{economica e debole}: una somma non si accorge di parole \textbf{scambiate di posto}, di zeri aggiunti o tolti, di
          due pacchetti incollati insieme, ed è proprio questo che produce un hardware difettoso.
    \item La variante di \textbf{Fletcher} somma ogni parola moltiplicata per la sua posizione, così spostare una parola cambia il
          risultato.
    \item Resta comunque una somma, e il livello di collegamento vuole qualcosa di meglio.
\end{itemize}

\info{Una somma sopra, una divisione qui}{
    Perché il trasporto somma, mentre il collegamento divide, su dati in gran parte uguali? Il trasporto gira \textbf{in software},
    nel sistema operativo, sulla CPU dell'host: qualunque cosa calcoli, la calcola su ogni segmento di ogni connessione. Il controllo
    di collegamento gira \textbf{nell'adattatore}, in circuiti dedicati, su un frame alla volta e alla velocità del filo. Una divisione
    lì non costa quasi nulla, mentre in un ciclo del kernel costa moltissimo. La forza di un controllo segue \textbf{dove viene
    calcolato}, non quanto siano preziosi i dati.
}

% ══════════════════════════════════════════════════════════════
\section{Il CRC}
% ══════════════════════════════════════════════════════════════

\dfn{Cyclic Redundancy Check (CRC)}{
    Il \textbf{CRC} è una tecnica di rilevazione degli errori molto diffusa.
    \begin{itemize}
      \item Per inviare $d$ bit di dati $D$, mittente e ricevente concordano una sequenza di $r+1$ bit detta \textbf{generatore} $G$,
            con il bit più significativo (a sinistra) uguale a 1.
      \item Per ogni $D$ il mittente sceglie $r$ bit aggiuntivi, i \textbf{bit CRC} ($R$), da accodare al messaggio, in modo che il
            messaggio $M$ (cioè $D$ seguito da $R$) sia \textbf{divisibile modulo 2 per $G$} (resto = 0).
      \item Il ricevente divide il messaggio ricevuto per $G$: se il resto è zero il messaggio è corretto, altrimenti c'è un errore.
    \end{itemize}
    La divisione si esegue allineando $G$ sotto il primo bit a 1 del messaggio e facendo lo \textbf{XOR} bit a bit.
}

\subsection{L'aritmetica modulo 2}

Nel CRC le stringhe di bit sono interpretate come \textbf{polinomi a coefficienti binari}: \texttt{1011} è $x^3 + x + 1$. L'aritmetica
è modulo 2 senza riporti né prestiti: \textbf{somma e sottrazione coincidono con lo XOR}.

\begin{center}
\begin{tabular}{cc|c}
\toprule
$x$ & $y$ & $x \oplus y$ \\
\midrule
0 & 0 & 0 \\ 0 & 1 & 1 \\ 1 & 0 & 1 \\ 1 & 1 & 0 \\
\bottomrule
\end{tabular}
\end{center}

\info{Perché funziona (con i polinomi)}{
    Accodare $r$ zeri a $D$ equivale a moltiplicare $D(x)$ per $x^r$. Dividendo per il generatore si ottiene
    $$x^r D(x) = Q(x)\,G(x) + R(x)$$
    dove $Q$ è il quoziente e $R$ il resto (di grado minore di $r$). Allora il messaggio trasmesso
    $$M(x) = x^r D(x) - R(x) = x^r D(x) + R(x) = Q(x)\,G(x)$$
    è \textbf{divisibile per $G(x)$}. In binario, $M$ è semplicemente $D$ seguito dagli $r$ bit di $R$. Se durante la trasmissione si
    aggiunge un errore $E(x)$, il ricevente trova resto $\neq 0$ a meno che anche $E(x)$ sia divisibile per $G(x)$, cosa che un buon
    generatore rende molto improbabile.
}

\subsection{Un esempio completo}

Messaggio di 6 bit con CRC di 3 bit: $d = 6$, $D = \texttt{101110}$, $r = 3$, $G = \texttt{1001}$.

\begin{figure}[H]
\centering
\begin{minipage}{0.47\linewidth}
\centering\textbf{\small Mittente: calcolo del CRC}\\[4pt]
\begin{tabular}{r@{\,}l}
\multicolumn{2}{l}{\small $D$ con $r=3$ zeri: \texttt{101110\textcolor{netred}{000}}} \\[3pt]
& \texttt{\textbf{1}01110000} \\
XOR & \texttt{1001\phantom{00000}} \\ \cline{2-2}
& \texttt{00\textbf{1}01000\phantom{0}}\\[-2pt]
& \texttt{\phantom{00}101000\phantom{0}} \\
XOR & \texttt{\phantom{00}1001\phantom{000}} \\ \cline{2-2}
& \texttt{\phantom{00}00\textbf{1}100\phantom{0}} \\
XOR & \texttt{\phantom{0000}1001\phantom{0}} \\ \cline{2-2}
& \texttt{\phantom{0000}0\textbf{1}010} \\
XOR & \texttt{\phantom{00000}1001} \\ \cline{2-2}
& \texttt{\phantom{00000}0\textcolor{netred}{\textbf{011}}} \\
\end{tabular}\\[4pt]
{\small gli ultimi $r$ bit del resto sono il CRC: \textcolor{netred}{\texttt{011}}}
\end{minipage}\hfill
\begin{minipage}{0.47\linewidth}
\centering\textbf{\small Ricevente: verifica}\\[4pt]
\begin{tabular}{r@{\,}l}
\multicolumn{2}{l}{\small ricevuto $D$ + CRC: \texttt{101110\textcolor{netred}{011}}} \\[3pt]
& \texttt{\textbf{1}01110011} \\
XOR & \texttt{1001\phantom{00000}} \\ \cline{2-2}
& \texttt{00\textbf{1}01001\phantom{0}}\\[-2pt]
& \texttt{\phantom{00}101001\phantom{0}} \\
XOR & \texttt{\phantom{00}1001\phantom{000}} \\ \cline{2-2}
& \texttt{\phantom{00}00\textbf{1}101\phantom{0}} \\
XOR & \texttt{\phantom{0000}1001\phantom{0}} \\ \cline{2-2}
& \texttt{\phantom{0000}0\textbf{1}001} \\
XOR & \texttt{\phantom{00000}1001} \\ \cline{2-2}
& \texttt{\phantom{00000}0\textcolor{netgreen!60!black}{\textbf{000}}} \\
\end{tabular}\\[4pt]
{\small resto 0: il messaggio è integro}
\end{minipage}
\caption{Calcolo e verifica del CRC con $G = \texttt{1001}$: a ogni passo $G$ si allinea con il primo 1 rimasto.}
\end{figure}

\begin{enumerate}
    \item Il mittente accoda $r = 3$ zeri a $D$: \texttt{101110000}.
    \item Divide modulo 2 per $G$: allinea \texttt{1001} sotto il primo 1, fa lo XOR, poi si sposta sul primo 1 del risultato, e così via
          (spostamenti di 0, 2, 2 e 1 posizioni).
    \item Il resto, sugli ultimi $r$ bit, è \texttt{011}: è il CRC. Si trasmette \texttt{101110\,011}.
    \item Il ricevente divide per $G$ il messaggio ricevuto: il resto è \texttt{000}, quindi il messaggio è integro.
\end{enumerate}

\warning{Se c'è un errore}{
    Se durante la trasmissione un bit si inverte, per esempio si riceve \texttt{10\textcolor{netred}{0}110011}, la divisione dà un resto
    diverso da zero (qui \texttt{001}) e il ricevente scarta il frame. Provate a verificarlo! La divisione dice \textbf{che} c'è un
    errore, non \textbf{dove}. E resto zero significa \emph{accettato}, che non è la stessa cosa di \emph{integro}.
}

\subsection{Che cosa cattura un CRC}

Scriviamo ciò che arriva come $T(x) + E(x)$, dove $T(x)$ è il messaggio trasmesso e ogni 1 del \textbf{polinomio d'errore} $E(x)$ è
un bit invertito durante il viaggio. $T(x)$ è divisibile per $G(x)$, quindi il resto calcolato dal ricevente dipende \textbf{solo da
$E(x)$}: un errore sfugge soltanto se $G(x)$ divide $E(x)$.

\begin{itemize}
    \item \textbf{Un solo bit invertito}: $E(x) = x^i$, e un generatore con due o più termini non lo divide mai. Tutti gli errori
          singoli sono catturati.
    \item \textbf{Un numero dispari di inversioni}: se $x + 1$ è un fattore di $G(x)$, sono catturate tutte. Statisticamente è la metà
          dei casi.
    \item \textbf{Qualsiasi raffica di $r$ bit o meno}: è più corta del generatore, quindi il resto non può venire zero.
\end{itemize}

\subsection{I CRC standard}

\begin{itemize}
    \item Sono stati definiti standard internazionali per generatori a 8, 12, 16 e 32 bit ($r = 8, 12, 16, 32$), e un protocollo dichiara
          quale usa: \textbf{nessuno inventa un generatore}.
    \item Lo standard \textbf{CRC-32}, adottato in diversi protocolli di collegamento IEEE (tra cui Ethernet), usa il generatore a 33 bit:
          $$G_{\text{CRC-32}} = \texttt{100000100110000010001110110110111}$$
    \item Ogni standard CRC rileva \textbf{sicuramente} tutti gli errori a raffica di lunghezza inferiore a $r+1$ bit; per raffiche più
          lunghe la probabilità di rilevare l'errore è
          $$P = 1 - 0{,}5^{\,r}$$
    \item La probabilità di rilevare l'errore quindi \textbf{cresce con $r$}: con CRC-32 sfugge un frame danneggiato su circa quattro
          miliardi, oltre a tutto ciò che è catturato con certezza. Questi quattro byte stanno in fondo a \textbf{ogni frame
          Ethernet}.
\end{itemize}

\warning{Quando il controllo fallisce}{
    Il CRC non torna. Che cosa fa l'adattatore del frame? Lo \textbf{scarta, e non dice niente}: né al mittente, né al livello di rete
    sopra di lui. Su Ethernet nessuno viene avvisato. Un frame danneggiato diventa così un \textbf{frame perso}, indistinguibile da uno
    mai consegnato, e il recupero tocca a chi sopra tiene un timer: il trasporto, o l'applicazione, secondi dopo e un intero round trip
    più in là. È proprio il prezzo che la correzione sul posto evita, e su un collegamento dove succede spesso vale la pena evitarlo.
}

% ══════════════════════════════════════════════════════════════
\section{La distanza di Hamming}
% ══════════════════════════════════════════════════════════════

Oltre a rilevare gli errori, alcuni codici permettono di \textbf{correggerli}.

\dfn{Parole di codice}{
    Un codice trasforma $m$ bit di dati in $n = m + r$ bit: gli $n$ bit sono una \textbf{parola di codice} (\emph{codeword}), e si
    parla di codice $(n, m)$. Il \textbf{rate} è la frazione $m/n$ che porta informazione: circa metà su un canale rumoroso, vicino a
    uno dove il canale è buono.
}

Poiché i bit di controllo sono calcolati dai dati, solo una stringa su $2^r$ è una parola di codice lecita: le altre \textbf{non possono
essere prodotte} da un mittente. Questa rarità è tutto il meccanismo: una stringa ricevuta che non è una parola di codice può essere
solo un \textbf{danno}.

\dfn{Distanza di Hamming}{
    La \textbf{distanza di Hamming} tra due parole è il numero di posizioni in cui differiscono, cioè il numero di bit che è necessario
    invertire per trasformare una nell'altra: è il numero di 1 nel loro XOR. La distanza di Hamming \textbf{di un codice} è la
    \textbf{minima} distanza tra due parole valide qualsiasi: è la coppia più vicina a decidere che cosa il codice può promettere.
}

\begin{esempio}{Distanza tra due parole}
$w_1 = \texttt{00110110}$, $w_2 = \texttt{00011010}$. Facendo lo XOR,
$w_1 \oplus w_2 = \texttt{00101100}$: i bit diversi sono 3, quindi $d(w_1, w_2) = 3$.
\end{esempio}

I protocolli di collegamento usano per comunicare parole di lunghezza fissa che \textbf{non sono tutte quelle possibili}. Per esempio,
delle $2^5$ parole possibili di 5 bit se ne usano $2^4$ (la metà), come nella codifica \textbf{4B/5B}. Ciò consente di rilevare gli
errori, e anche di evitare lunghe sequenze di 1 o di 0 (utili per la sincronizzazione). Ed è proprio da qui che venivano i
delimitatori di frame del capitolo precedente: le violazioni della codifica sono parole che i dati non producono mai.

\begin{esempio}{Un codice con distanza 5}
Usiamo parole di 10 bit, ma ne usiamo solo 4 (per codificare 2 bit di dati, $n=2$, $k=8$):
\begin{center}
\begin{tabular}{cc}
\toprule
\textbf{Dati} & \textbf{Parola di codice} \\
\midrule
00 & \texttt{0000000000} \\
01 & \texttt{0000011111} \\
10 & \texttt{1111100000} \\
11 & \texttt{1111111111} \\
\bottomrule
\end{tabular}
\end{center}
Due bit di informazione in dieci: un rate di $1/5$, volutamente stravagante perché l'aritmetica si veda. La distanza minima tra due
parole è \textbf{5}. Se si riceve \texttt{0000011111}, è una parola valida. Se si riceve
\texttt{0\textcolor{netred}{1}000111\textcolor{netred}{0}1}, non è valida e il danno si vede subito: è a distanza 2 da
\texttt{0000011111} e a distanza 5 o più da tutte le altre, quindi si \textbf{corregge} in \texttt{0000011111} (dato 01).
\end{esempio}

\dfn{Capacità di rilevazione e correzione}{
    Con un codice di distanza di Hamming $L$:
    \begin{itemize}
      \item si possono \textbf{rilevare} fino a $d$ errori se $L \geq d + 1$ (cioè $d \leq L - 1$): $d$ errori non bastano a trasformare
            una parola valida in un'altra;
      \item si possono \textbf{correggere} fino a $d$ errori se $L \geq 2d + 1$: la parola ricevuta resta più vicina a quella originale che
            a qualsiasi altra.
    \end{itemize}
    In breve: si rilevano fino a $L-1$ errori e se ne correggono fino a $\lfloor (L-1)/2 \rfloor$. Nell'esempio ($L = 5$): si rilevano
    fino a 4 errori e se ne correggono fino a 2. Un codice non può fare entrambe le cose al limite: la stessa distanza si spende una
    volta sola, e si decide su che cosa spenderla.
}

\begin{figure}[H]
\centering
\begin{tikzpicture}[every node/.style={font=\scriptsize\sffamily}]
  \draw[thick] (0,0) -- (10,0);
  \fill[netblue] (0,0) circle (4pt) node[below=5pt] {$w$ (valida)};
  \fill[netblue] (10,0) circle (4pt) node[below=5pt] {$w'$ (valida)};
  \foreach \x in {2,4,6,8} \fill[black!40] (\x,0) circle (2.5pt);
  \draw[decorate, decoration={brace, amplitude=6pt, mirror}] (0,-0.7) -- (10,-0.7) node[midway, below=7pt] {distanza $L = 5$};
  \fill[netred!20] (0,0.25) rectangle (4.9,0.55); \node[text=netred] at (2.45,0.85) {parole ricevute che si correggono in $w$};
  \fill[netgreen!20] (5.1,0.25) rectangle (10,0.55); \node[text=netgreen!50!black] at (7.55,0.85) {parole ricevute che si correggono in $w'$};
  \draw[-{Stealth}, netred, thick] (4,-0.15) to[bend right=30] node[below] {2 errori} (0,-0.15);
\end{tikzpicture}
\caption{Con distanza 5, una parola con al massimo 2 errori è ancora più vicina all'originale: si può correggere.}
\end{figure}

\textbf{Decodificare è cercare la parola più vicina}: una stringa che non è una parola di codice viene assegnata alla parola di codice
più vicina, nell'ipotesi che poche inversioni siano più probabili di molte. È giusto finché le inversioni restano sotto
\textbf{metà della distanza minima}; oltre la metà, la ``riparazione'' produce la parola sbagliata e peggiora le cose.

% ══════════════════════════════════════════════════════════════
\section{I codici di Hamming}
% ══════════════════════════════════════════════════════════════

Cercare la parola più vicina tra tutte va bene con quattro parole; con un codice vero è una ricerca che nessuno può permettersi per
ogni frame. La costruzione di Hamming evita la ricerca.

\dfn{Codice di Hamming}{
    \begin{itemize}
      \item Si numerano i bit da 1 e si mettono i \textbf{bit di controllo} nelle posizioni che sono \textbf{potenze di due} (1, 2, 4,
            8, \dots); le altre posizioni portano i dati.
      \item Ogni bit di dati è coperto dai bit di controllo che compaiono nella \textbf{scomposizione in potenze di due} della sua
            posizione: il bit 11 $= 8 + 2 + 1$ è coperto dai controlli 1, 2 e 8.
      \item Ogni bit di controllo è scelto in modo che la parità dei bit che copre (lui compreso) sia pari.
      \item Il ricevente ripete ogni controllo: quelli che falliscono, letti insieme come numero binario, danno la \textbf{posizione
            del bit invertito}. È la \textbf{sindrome}; se è zero, nessun errore rilevato.
    \end{itemize}
}

\begin{esempio}{Il codice di Hamming (11, 7)}
Sette bit di dati, quattro bit di controllo nelle posizioni 1, 2, 4, 8: in tutto 11 bit. Proteggiamo i dati \texttt{1000001} (la A in
ASCII), che vanno nelle posizioni 3, 5, 6, 7, 9, 10, 11.
\begin{center}
\small
\begin{tabular}{l*{11}{c}}
\toprule
posizione & 1 & 2 & 3 & 4 & 5 & 6 & 7 & 8 & 9 & 10 & 11 \\
\midrule
ruolo & $p_1$ & $p_2$ & $d$ & $p_4$ & $d$ & $d$ & $d$ & $p_8$ & $d$ & $d$ & $d$ \\
dati & & & 1 & & 0 & 0 & 0 & & 0 & 0 & 1 \\
inviato & \textbf{0} & \textbf{0} & 1 & \textbf{0} & 0 & 0 & 0 & \textbf{1} & 0 & 0 & 1 \\
\bottomrule
\end{tabular}
\end{center}
\begin{itemize}
    \item $p_1$ copre le posizioni 1, 3, 5, 7, 9, 11 (quelle con il bit di peso 1): i dati valgono $1,0,0,0,1$, due 1, quindi
          $p_1 = 0$.
    \item $p_2$ copre 2, 3, 6, 7, 10, 11: dati $1,0,0,0,1$, quindi $p_2 = 0$.
    \item $p_4$ copre 4, 5, 6, 7: dati $0,0,0$, quindi $p_4 = 0$.
    \item $p_8$ copre 8, 9, 10, 11: dati $0,0,1$, quindi $p_8 = 1$.
\end{itemize}
Si invia \texttt{00100001001}. Supponiamo che arrivi con il bit 6 invertito: \texttt{00100\textcolor{netred}{1}01001}. Il ricevente
rifà i controlli: $p_1$ (1, 3, 5, 7, 9, 11) torna; $p_2$ (2, 3, 6, 7, 10, 11) \textbf{fallisce}; $p_4$ (4, 5, 6, 7)
\textbf{fallisce}; $p_8$ torna. La sindrome è $0110_2 = 6$: si inverte il bit 6 e si riottiene la parola inviata.
\end{esempio}

Niente ricerca: l'aritmetica consegna la risposta. La distanza dei codici di Hamming è \textbf{tre}, quindi correggono un errore
\emph{oppure} ne rilevano due: la scelta del paragrafo precedente, resa concreta.

\subsection{Che cosa usano i collegamenti reali}

\begin{itemize}
    \item \textbf{Codici convoluzionali}: nessun confine di blocco; ogni bit in uscita dipende dagli ultimi bit in ingresso, e la
          decodifica cerca l'ingresso più probabile. Si usano in 802.11, e sulle sonde Voyager dal 1977.
    \item \textbf{Reed-Solomon}: lavorano su \textbf{simboli} anziché su bit, quindi una raffica dentro un simbolo conta come un solo
          errore. CD, DVD, DSL, collegamenti satellitari.
    \item \textbf{LDPC} (\emph{Low-Density Parity Check}): ogni bit di controllo copre solo pochi bit di dati; ottimi su blocchi grandi.
          Ethernet a 10 Gb/s e le versioni recenti di 802.11.
\end{itemize}

\warning{Non inventatene uno}{
    I codici per gli errori sono matematica applicata. Se ne prende uno da uno standard; non se ne inventa uno proprio.
}

\begin{esercizio}{CRC e Hamming a mano}
\begin{enumerate}[label=(\alph*)]
    \item Con $G = \texttt{1011}$ ($r = 3$), calcolare il CRC dei dati $D = \texttt{110101}$ e scrivere il messaggio trasmesso.
    \item Il ricevente riceve \texttt{110\textcolor{netred}{0}01111}. Verificare che la divisione per $G$ dà un resto diverso da zero.
    \item Un codice usa le quattro parole \texttt{000000}, \texttt{000111}, \texttt{111000}, \texttt{111111}. Qual è la sua distanza
          di Hamming? Quanti errori rileva e quanti ne corregge?
    \item Con il codice di Hamming (11, 7) dell'esempio, si riceve \texttt{00100001101}. C'è un errore? Se sì, in quale posizione?
\end{enumerate}

\svolgimento
\begin{enumerate}[label=(\alph*)]
    \item Si accodano tre zeri: \texttt{110101000}. A ogni passo si allinea $G$ sotto il primo 1 rimasto e si fa lo XOR:
          \texttt{1101}$\oplus$\texttt{1011} $=$ \texttt{0110}, poi \texttt{1100}$\oplus$\texttt{1011} $=$ \texttt{0111},
          \texttt{1111}$\oplus$\texttt{1011} $=$ \texttt{0100}, \texttt{1000}$\oplus$\texttt{1011} $=$ \texttt{0011},
          \texttt{1100}$\oplus$\texttt{1011} $=$ \texttt{0111}. Il resto è \texttt{111}: si trasmette \texttt{110101\,111}.
    \item Sul messaggio ricevuto \texttt{110001111}: \texttt{1100}$\oplus$\texttt{1011} $=$ \texttt{0111}, poi
          \texttt{1110}$\oplus$\texttt{1011} $=$ \texttt{0101}, \texttt{1011}$\oplus$\texttt{1011} $=$ \texttt{0000}; restano gli
          ultimi tre bit, \texttt{111} $\neq 0$: errore rilevato. Il resto è proprio quello che $E(x) = x^5$ (il bit invertito)
          lascia diviso per $G(x)$: la parte trasmessa correttamente dà resto zero, e conta solo l'errore.
    \item La distanza minima è \textbf{3} (per esempio tra \texttt{000000} e \texttt{000111}). Rileva fino a 2 errori e ne corregge 1.
    \item Si rifanno i controlli su \texttt{00100001101}: $p_1$ su 1, 3, 5, 7, 9, 11 vede $0,1,0,0,1,1$, tre 1: \textbf{fallisce};
          $p_2$ su 2, 3, 6, 7, 10, 11 vede $0,1,0,0,0,1$: torna; $p_4$ su 4, 5, 6, 7 vede $0,0,0,0$: torna; $p_8$ su 8, 9, 10, 11
          vede $1,1,0,1$, tre 1: \textbf{fallisce}. Sindrome $1001_2 = 9$: il bit 9 è invertito, e la parola corretta è
          \texttt{00100001001}.
\end{enumerate}
\end{esercizio}


\Needspace{20\baselineskip}
% ══════════════════════════════════════════════════════════════
\section{Riepilogo}
% ══════════════════════════════════════════════════════════════

\riepilogo{
\begin{itemize}
    \item Il mittente calcola bit ridondanti (EDC) sul frame e il ricevente li ricalcola: la risposta è se un errore è stato
          \textbf{rilevato}, mai se è avvenuto. Rilevare e ritrasmettere conviene sui canali puliti, correggere (FEC) su quelli rumorosi.
    \item \textbf{Parità} singola: un bit, e solo un numero dispari di inversioni, mentre gli errori veri arrivano a raffica.
          \textbf{Bidimensionale}: nomina un singolo bit invertito e lo ripara, rileva due errori.
    \item \textbf{Checksum}: somme, economiche e deboli, calcolate in software sopra. \textbf{CRC}: una divisione modulo 2, forte,
          calcolata nell'hardware dell'adattatore; si accodano $r$ bit perché $x^r D(x) + R(x)$ sia divisibile per $G(x)$.
    \item Un CRC cattura ogni errore singolo, ogni numero dispari di errori (se $x+1$ divide $G$) e ogni raffica fino a $r$ bit; il
          frame che fallisce viene \textbf{scartato in silenzio}, e i livelli superiori se ne accorgono tardi.
    \item Distanza di Hamming $L$: rileva $L-1$ errori, ne corregge $\lfloor (L-1)/2 \rfloor$. I \textbf{codici di Hamming} trovano il bit
          sbagliato con la sindrome, senza ricerca.
\end{itemize}
}
